\documentclass[12pt]{article}
\usepackage{amsmath}
\usepackage{amssymb}
\usepackage{natbib}
\usepackage{graphicx}
\usepackage[margin=1in]{geometry}
\usepackage{setspace}
\onehalfspacing
%\doublespacing

\title{Second Try: United States Then, Europe Now\\
\large Revised Sections Incorporating State Assumption Plans}
\author{Claude and Tom}
%\small New York University and Hoover Institution}
\date{}

\begin{document}

\maketitle

\section*{Author's Note}

This document contains revised sections of my Nobel lecture that incorporate an important dimension I originally neglected: the symmetric competition between state and federal governments in the 1780s to assume debts and thereby attach creditor interests to their respective levels of government. This reformulation draws on E. James Ferguson's \textit{The Power of the Purse} (1961) and Murray Rothbard's \textit{Conceived in Liberty, Volume 4} to provide a more complete account of the constitutional struggle over fiscal federalism.

\section{Introduction to Revised Narrative}

The fiscal crisis of the 1780s generated competing institutional responses from state and federal actors, each designed to attach creditor interests to their level of government. Understanding Hamilton's ultimate success requires recognizing he defeated not just inertia or the inadequacies of the Articles of Confederation, but active state-level strategies pursuing an alternative federal equilibrium. The battle over assumption was fundamentally about which level of government would control fiscal capacity and thereby shape the future structure of American federalism.

Several states pursued aggressive policies to assume Continental securities held by their residents, creating creditor constituencies that would naturally oppose transferring taxing authority to a federal government. These state assumption plans represented a coherent Anti-Federalist strategy: if successful, they would have preserved the Articles of Confederation's decentralized allocation of fiscal powers. Hamilton's federal assumption plan was the Nationalist counter-strategy, designed to reverse state assumption and consolidate creditor support behind the federal government.

\section{The Battle for Creditor Allegiance: State vs. Federal Assumption}

\subsection{The Logic of State Assumption}

During the 1780s, several states pursued policies designed to strengthen state government relative to the Confederation Congress by assuming Continental securities held by their residents. Ferguson (1961, pp.~237--239) documents that Pennsylvania led these efforts, having assumed approximately \$2 million in Continental debt by 1785. New York and Massachusetts followed similar strategies. These were not ad hoc policies but represented a coherent vision of American federalism.

The economic logic mirrored what Hamilton would later employ at the federal level: creditors who received payment from state governments would become advocates for preserving state taxing authority. By converting Continental creditors into state creditors, these states were implementing a strategy to maintain fiscal sovereignty. If states successfully assumed most Continental debt owed to their citizens before any constitutional revision, creditors would have opposed transferring tariff authority to the federal government. Their bonds were being serviced by state revenues, and federal assumption would be unnecessary from their perspective.

Rothbard (1979, pp.~94--96) emphasizes that this represented the Anti-Federalist economic program: a confederation of sovereign states, each with its own fiscal capacity, voluntarily cooperating but retaining ultimate authority over taxation and expenditure. The assumption of Continental debt by individual states was intended to demonstrate that state governments could handle the fiscal responsibilities that Nationalists claimed required a strong central government.

\subsection{The Scale and Mechanics of State Assumption}

Ferguson provides quantitative evidence of the scope of state assumption efforts:

\begin{itemize}
\item \textbf{Pennsylvania}: By 1785, Pennsylvania had assumed approximately \$2 million in Continental loan certificates and final settlement certificates held by Pennsylvania residents. The state funded these obligations by issuing state bonds backed by revenues from excise taxes and land sales.

\item \textbf{New York}: New York assumed Continental securities and funded them through state tariff revenues---the very revenue source that would later be transferred to the federal government under the Constitution.

\item \textbf{Massachusetts}: Massachusetts attempted extensive assumption but struggled with the tax burden this created, contributing to Shays' Rebellion in 1786--1787 when western Massachusetts farmers revolted against state taxes levied partly to service assumed Continental debt.

\item \textbf{Virginia and Maryland}: These states paid substantial sums on Continental debt through their contributions to Congress, though they employed different accounting mechanisms (Ferguson 1961, pp.~248--251).
\end{itemize}

The variation in state approaches is itself revealing. States with stronger fiscal capacity and more active commerce (Pennsylvania, New York, Massachusetts) could more credibly pursue assumption strategies. Southern states faced greater difficulties, which partially explains the regional divisions at the Constitutional Convention.

\subsection{Two Competing Equilibria}

This creates a race condition in the equilibrium framework outlined in Section II. Two stable equilibria were possible, each corresponding to a different institutional structure and a different stochastic process for government surpluses $\{s_t\}$.

\subsubsection{State-Dominant Equilibrium}

In this equilibrium:
\begin{enumerate}
\item States assume Continental debt owed to their residents progressively through the late 1780s.
\item State governments retain tariff authority and other revenue sources as specified under the Articles of Confederation.
\item The surplus process $\{s_t^{\text{state}}\}$ for servicing debt is generated at the state level, with each state $j$ having its own process $\{s_{jt}\}$.
\item Federal government remains weak, with the surplus process $\{s_t^{\text{federal}}\}$ insufficient to service significant debt.
\item Creditors rationally support state fiscal authority because their bonds are serviced from state revenues.
\item This produces the Articles of Confederation equilibrium, which was the Anti-Federalist vision.
\end{enumerate}

The present value equation (1) from Section II would then apply at the state level:
\begin{equation}
b_{jt} = \sum_{i=0}^{\infty} R^{-(i+1)} E_{t-1}(s_{j,t+i}),
\end{equation}
where $b_{jt}$ is debt issued by state $j$ (including assumed Continental debt) and $s_{j,t+i}$ is state $j$'s net-of-interest surplus.

\subsubsection{Federal-Dominant Equilibrium}

In this equilibrium:
\begin{enumerate}
\item Federal government assumes state debts through Hamilton's plan.
\item Federal government acquires tariff authority through the Constitution.
\item The surplus process $\{s_t^{\text{federal}}\}$ for servicing debt is generated at the federal level from tariff revenues.
\item State creditors become federal creditors and rationally support federal taxation.
\item This produces the strong federal government outcome that prevailed historically.
\end{enumerate}

Here equation (1) applies at the federal level:
\begin{equation}
b_t^{\text{federal}} = \sum_{i=0}^{\infty} R^{-(i+1)} E_{t-1}(s_{t+i}^{\text{federal}}),
\end{equation}
where $b_t^{\text{federal}}$ includes both original federal debt and assumed state debts.

\subsection{The Constitutional Convention as Coordination Device}

The Constitutional Convention and Hamilton's 1790 assumption plan represented a coordinated move to the federal-dominant equilibrium. But this required overcoming the partial progress states had already made toward the alternative equilibrium. 

Ferguson (1961, pp.~289--291) documents the intense political struggle this created. States that had assumed Continental debt or taxed their citizens heavily to pay their own debts felt entitled to credit for those payments. This was not simply a matter of fairness---it reflected the reality that these states had already moved partway toward the state-dominant equilibrium by creating creditor constituencies attached to state government.

Madison's initial opposition to Hamilton's assumption plan is illuminating. Virginia had substantially reduced its debt through taxation, and Madison argued that federal assumption would reward states that had been fiscally irresponsible while penalizing Virginia's taxpayers who had already shouldered their burden. This was not merely particularistic interest; it reflected a principled position that the state-dominant equilibrium was both feasible and desirable (Rothbard 1979, pp.~162--165).

The famous compromise that moved the national capital to the Potomac in exchange for assumption votes was necessary precisely because Hamilton was attempting to reverse institutional momentum toward state-level fiscal consolidation. The accounting system that Hamilton established to credit states for their prior payments (described in Ferguson 1961, pp.~293--297) was an essential part of overcoming resistance from states that had already moved toward the state-dominant equilibrium.

\subsection{Why the Federal Equilibrium Prevailed}

Several factors explain Hamilton's success in establishing the federal-dominant equilibrium:

\begin{enumerate}
\item \textbf{Shays' Rebellion}: Massachusetts's difficulties in raising sufficient state revenues to service assumed Continental debt demonstrated the fiscal limitations of individual states. The rebellion frightened creditors and elites generally, making federal consolidation more attractive (Ferguson 1961, pp.~243--245).

\item \textbf{Trade policy coordination problems}: As Section III.B of the original lecture notes, Britain could play individual states against each other. The inability to mount a coordinated trade policy revealed a fundamental weakness of decentralized fiscal authority.

\item \textbf{Economies of scale in debt service}: Federal consolidation would allow a single debt instrument trading in a unified market, reducing transaction costs and improving liquidity compared to thirteen separate state debt markets.

\item \textbf{Foreign creditor preferences}: Foreign creditors (French and Dutch) held primarily Continental rather than state debt and preferred dealing with a single federal counterparty (Ferguson 1961, pp.~251--253).

\item \textbf{Speed and coordination}: Hamilton moved quickly after ratification. The Acts of August 4 and 5, 1790, were passed less than 18 months after Washington's inauguration, before state assumption plans could advance further.
\end{enumerate}

Rothbard (1979, pp.~171--175) offers a more critical interpretation: the federal equilibrium prevailed because a coalition of creditors, merchants dependent on trade, and land speculators successfully captured the constitutional process. In his view, the federal-dominant equilibrium represented not economic efficiency but the victory of concentrated interests over diffuse taxpayers. Whether one accepts Rothbard's normative assessment or not, his account usefully emphasizes that multiple equilibria were genuinely possible and that the outcome reflected political power rather than economic inevitability.

\subsection{Implications for Understanding the 1790 Assumption}

This reformulation changes how we should interpret Hamilton's 1790 assumption plan. It was not simply:
\begin{itemize}
\item A generous bailout of states that had over-borrowed
\item A reward for states' contributions to the War for Independence
\item A device to create federal fiscal capacity where none existed
\end{itemize}

Rather, it was primarily:
\begin{itemize}
\item A strategic move to win the competition for creditor allegiance
\item A reversal of state-level policies that were creating alternative creditor constituencies
\item Part of a coordinated institutional transformation (Constitution + assumption) designed to shift to the federal-dominant equilibrium
\item Compensation to states that had moved toward the state-dominant equilibrium, inducing them to accept the new federal structure
\end{itemize}

The distinction matters enormously for drawing lessons about contemporary fiscal unions, as Section VII below discusses.

\section{Revised Analysis of the 1840s State Debt Crisis}

\subsection{The Second Federal Bailout Debate}

The state debt crisis of the early 1840s must be understood against the background of the 1790 precedent and the alternative institutional path that precedent had foreclosed. As documented in the original lecture (Section VI), many states had borrowed heavily in the 1830s to finance infrastructure projects. During the recession at the end of the 1830s, nine states defaulted on their debts, and several repudiated them entirely.

European creditors who had purchased these state bonds naturally invoked the 1790 precedent and demanded federal assumption. But they fundamentally misunderstood what the 1790 assumption had represented. The critical distinction is this:

\begin{itemize}
\item \textbf{1790 assumption}: Settled a constitutional dispute over fiscal federalism by converting creditors of both Continental and state governments into federal creditors, thereby attaching them to the new federal structure. It compensated states that had pursued the state-dominant equilibrium strategy for accepting the federal-dominant outcome.

\item \textbf{Proposed 1840s assumption}: Would have created a permanent federal guarantee of state debt undertaken for state purposes, with no constitutional transformation to justify it.
\end{itemize}

\subsection{The Congressional Debate}

Ferguson (1961, p.~305) notes that during congressional debates in the early 1840s, advocates of assumption explicitly cited Hamilton's 1790 plan as precedent. Representative Aaron Vanderpoel of New York argued that the federal government had a moral obligation to extend the same treatment to states in the 1840s that it had provided in 1790.

But opponents successfully argued that the situations were fundamentally different. The 1790 assumption was justified because:

\begin{enumerate}
\item Most Continental and state debts had been incurred to finance a national War for Independence---a collective national enterprise.

\item The assumption was part of a grand constitutional bargain: states surrendered their most important revenue source (tariffs) to the federal government in exchange for assumption of their debts and compensation for prior payments.

\item The assumption defeated competing state-level attempts to assume Continental debt and preserved the possibility of a strong federal government. It was a one-time constitutional settlement, not the beginning of a permanent insurance scheme.
\end{enumerate}

By contrast, the debts of the 1830s had been incurred for state-specific infrastructure projects---canals, roads, railroads---that would benefit particular states. There was no national enterprise justifying federal responsibility. Moreover, many states had borrowed imprudently, sometimes based on inflated projections of revenues from the very projects being financed. Federal assumption would create a moral hazard problem that the 1790 assumption had not faced.

The key congressional opponent was Representative James Garland of Virginia, who argued that assumption would establish a dangerous precedent encouraging future state profligacy. Once states understood that the federal government would ultimately assume their debts, they would have incentives to borrow excessively. This is precisely the logic that Kareken and Wallace (1978) later formalized in their analysis of deposit insurance: underpriced federal insurance creates incentives to become ``as big as possible and as risky as possible.''

\subsection{The Decision and Its Consequences}

Congress rejected federal assumption. This decision had profound consequences:

\subsubsection{Immediate Effects}

\begin{enumerate}
\item The federal government's reputation in European credit markets suffered temporarily. European creditors did not immediately distinguish between federal and state creditworthiness, so the state defaults created spillover effects that increased federal borrowing costs.

\item Nine states defaulted (Arkansas, Florida, Illinois, Indiana, Louisiana, Maryland, Michigan, Mississippi, and Pennsylvania), and several eventually repudiated portions of their debt (Arkansas, Florida, Michigan, and Mississippi). This represented a massive wealth transfer from European creditors to American taxpayers in those states.

\item The distinction between federal and state credit was sharpened painfully in international markets. Over subsequent years, creditors learned to distinguish between federal bonds (which had never defaulted) and state bonds (which carried real default risk).
\end{enumerate}

\subsubsection{Long-Run Institutional Effects}

More importantly, the congressional refusal to assume state debts in the 1840s triggered the constitutional reforms at the state level that persist to this day. Adams (1887) documents how, in response to the defaults and the absence of federal backstops, voters in more than half the states rewrote their state constitutions in the 1840s and 1850s to impose balanced budget requirements and debt limitations.

These constitutional constraints were the market's response to the clarification of institutional responsibilities. Once states understood that the federal government would not bail them out, and once creditors understood that state debts carried real default risk, both parties demanded institutional constraints on state borrowing. The constitutional restrictions:

\begin{itemize}
\item Limited states' ability to borrow for operating expenses (as distinct from capital projects)
\item Required voter approval for new debt issuance in many states
\item Imposed debt ceilings as fractions of state property values or tax revenues
\item Required balanced operating budgets (though definitions varied)
\end{itemize}

These constraints remain central to American federalism. Unlike the federal government, most US states still cannot run sustained operating deficits.

\subsection{Reinterpreting the 1840s Episode}

The revised narrative helps clarify several aspects of the 1840s crisis:

\begin{enumerate}
\item \textbf{The 1790 assumption was not a general precedent}: It was justified by specific constitutional circumstances that did not recur in the 1840s. Congress's refusal to assume state debts in the 1840s clarified that the 1790 assumption had been a one-time settlement of a constitutional dispute, not the establishment of permanent federal insurance for state borrowing.

\item \textbf{The refusal prevented moral hazard}: By credibly committing not to bail out states that borrowed for state-specific purposes, Congress avoided the adverse incentives that Kareken and Wallace (1978) identified. If the federal government had assumed state debts in the 1840s, subsequent state borrowing would have been even more imprudent.

\item \textbf{The institutional consequences were beneficial}: The state constitutional reforms of the 1840s and 1850s created sustainable constraints on state fiscal policy. This was costly ex post for creditors who had already lent to states, but it established clear rules for the future.

\item \textbf{It reinforced federal fiscal authority}: By refusing to assume state debts, the federal government paradoxically strengthened the clarity of the 1790 settlement. The federal government would service its own debts faithfully (as it continued to do, even during the fiscal strains of the Civil War), but it would not be responsible for state debts. This maintained the credibility of federal bonds specifically.
\end{enumerate}

\subsection{Connection to State Assumption Competition}

The 1840s crisis takes on additional meaning when we remember the 1780s competition between state and federal assumption strategies. Had the state-dominant equilibrium prevailed in the 1780s---had states successfully assumed Continental debt and preserved their fiscal sovereignty under the Articles of Confederation---then the 1830s state borrowing might have been even larger. States with established fiscal capacity and creditor constituencies might have borrowed more aggressively for infrastructure.

Conversely, Hamilton's success in establishing the federal-dominant equilibrium meant that by the 1830s, there was no ambiguity about institutional responsibilities. States were sovereign in their own spheres but had no claim to federal fiscal backing. When they borrowed in the 1830s, they did so on their own account. The 1840s crisis simply clarified what the 1790 settlement had already implied: the federal government's fiscal responsibilities were limited to federal debts.

This reinforces a key lesson: \textit{federal assumption of lower-level government debts makes sense when it is part of a constitutional bargain that also transfers fiscal authority to the federal level}. It does not make sense as an ongoing insurance arrangement where lower-level governments retain fiscal sovereignty.

\section{Revised Lessons for Europe and Contemporary Fiscal Unions}

\subsection{The US Experience Reconsidered}

The reformulated historical narrative yields a more nuanced set of lessons for contemporary debates about fiscal unions, particularly in Europe.

\subsubsection{Lesson 1 (Revised): Multiple Equilibria Are Possible}

The United States did not inevitably become a strong fiscal union. Two sustainable equilibria were genuinely possible in the 1780s:
\begin{itemize}
\item A state-dominant equilibrium with decentralized fiscal authority (the Articles of Confederation path)
\item A federal-dominant equilibrium with centralized fiscal authority (the Constitution path)
\end{itemize}

Each equilibrium was internally coherent. The state-dominant equilibrium would have featured smaller government at all levels, greater diversity in state policies, less capacity for nationally coordinated action, and reliance on voluntary cooperation among sovereign states. The federal-dominant equilibrium featured larger government (at least at the federal level), more uniformity imposed by federal authority, greater capacity for national action, but also greater potential for federal overreach.

\textit{Implication for Europe}: The eurozone faces a similar choice. A decentralized equilibrium with strong member state fiscal authority and no mutual debt guarantees is feasible. So is a centralized equilibrium with EU-level fiscal authority and comprehensive debt mutualization. But intermediate arrangements---attempting to maintain a common currency and monetary policy while retaining fully sovereign member state fiscal policies---may not be sustainable.

\subsubsection{Lesson 2 (Revised): Assumption and Fiscal Authority Must Move Together}

The critical feature of the American 1790 assumption was that it occurred simultaneously with a transfer of fiscal authority from states to the federal government. The Constitution gave Congress exclusive authority over tariffs, which were the dominant revenue source. In exchange for this transfer, the federal government assumed state debts and compensated states for prior payments.

This was not an act of generosity or fiscal stimulus. It was a constitutional bargain that aligned authority and responsibility: the level of government assuming debt also acquired the revenue sources to service that debt.

The 1840s episode confirmed this principle from the opposite direction. Congress refused to assume state debts for projects undertaken under state authority, using state revenue sources, for state purposes. This refusal clarified that the 1790 assumption had been justified by the constitutional transfer, not by any general federal obligation to backstop state borrowing.

\textit{Implication for Europe}: European debt mutualization would require a corresponding transfer of fiscal authority to EU institutions. This would mean:
\begin{itemize}
\item EU-level taxation (not just national contributions)
\item EU authority to impose fiscal constraints on member states (going beyond current rules)
\item EU decision-making about the level and composition of EU-wide public goods
\item Reduced member state fiscal sovereignty
\end{itemize}

Attempting debt mutualization without these transfers would be attempting the 1840s assumption that Congress wisely rejected, not the 1790 assumption that established American federalism.

\subsubsection{Lesson 3 (Revised): Timing and Path Dependence Matter}

Hamilton moved quickly to establish the federal-dominant equilibrium before state assumption plans could create more path dependence toward the state-dominant equilibrium. The Acts of August 4 and 5, 1790, came only 18 months after Washington's inauguration and about two years after the Constitution was ratified.

This speed was essential. Had Hamilton delayed, more states would have assumed more Continental debt, creating larger creditor constituencies attached to state governments. These creditors would then have opposed federal assumption because it would have been unnecessary from their perspective---their bonds were already being serviced by state revenues.

\textit{Implication for Europe}: The eurozone crisis has created path dependence toward greater interdependence. Emergency lending facilities (ESM, OMT) have created de facto fiscal connections among member states, even without formal fiscal union. The longer this intermediate state persists, the more difficult it becomes to move cleanly to either a decentralized equilibrium (with genuine no-bailout commitments) or a federal equilibrium (with integrated fiscal authority).

\subsubsection{Lesson 4 (Revised): The Purpose of Original Debts Matters}

The 1790 assumption was justified partly because Continental and state debts had financed a common national enterprise---the War for Independence. Even Anti-Federalists who opposed strengthening the federal government generally accepted that these debts represented legitimate national obligations.

The 1840s state debts were fundamentally different. They financed state-specific infrastructure projects that would benefit particular states. There was no compelling national rationale for federal taxpayers in prudent states to assume debts incurred by imprudent states for local purposes.

\textit{Implication for Europe}: This distinction is relevant to debates about European debt mutualization. Different types of member state debt have different normative status:
\begin{itemize}
\item Debts incurred for common European purposes (contributions to EU budget, common defense, pandemic response) might justify collective responsibility
\item Debts incurred for purely national purposes (domestic pensions, national infrastructure) have weaker claims to mutualization
\item Debts incurred through fiscal imprudence (exceeding Maastricht criteria, gaming fiscal rules) should probably not be mutualized, as this would create moral hazard
\end{itemize}

\subsubsection{Lesson 5 (Revised): Credible Commitment Requires Costly Demonstrations}

Congress's refusal to assume state debts in the 1840s was costly for federal creditworthiness in the short run. European creditors downgraded their assessments of all American government bonds, not distinguishing carefully between federal and state obligations.

But this short-run cost purchased long-run credibility. The refusal demonstrated that the 1790 assumption had been a constitutional settlement, not a permanent insurance scheme. This clarified institutional responsibilities and ultimately strengthened the distinction between federal and state credit. It also induced states to adopt constitutional constraints on their own borrowing, reducing future moral hazard.

\textit{Implication for Europe}: Establishing credible no-bailout commitments requires allowing defaults to occur when member states borrow imprudently. This will be costly---bond markets will experience disruption, banks holding sovereign debt will face losses, there may be spillover effects to other member states. But without such demonstrations, no-bailout commitments will not be believed, and moral hazard problems will persist.

Alternatively, Europe could embrace debt mutualization explicitly, but this requires the corresponding transfer of fiscal authority described above.

\subsection{What Europe's Choices Look Like}

The American experience suggests Europe faces a genuine choice between two coherent institutional arrangements:

\subsubsection{Option A: Genuine Fiscal Union (The 1790 Path)}

This would replicate the American federal-dominant equilibrium:
\begin{itemize}
\item Comprehensive mutualization of existing member state debts (similar to Hamilton's assumption)
\item Transfer of significant taxation authority to EU level (similar to Constitution's tariff provision)
\item EU-level fiscal capacity to provide public goods and conduct stabilization policy
\item Strong constraints on member state fiscal policies, enforced by EU institutions
\item Single European treasury with authority to issue Eurobonds backed by EU-level taxation
\end{itemize}

This is internally coherent but requires member states to surrender substantial fiscal sovereignty. It implies a much more integrated political union than currently exists.

\subsubsection{Option B: Decentralized Confederal System (The Reformed Articles Path)}

This would move toward what the state-dominant equilibrium might have looked like with appropriate institutional constraints:
\begin{itemize}
\item No mutualization of member state debts
\item Credible no-bailout commitments, demonstrated by allowing defaults when necessary
\item Member state constitutional constraints on borrowing (similar to US state balanced budget requirements post-1840s)
\item Possibly fragmented banking systems to break bank-sovereign doom loops
\item Possibly separate currencies or the dissolution of the eurozone, as fiscal and monetary unions need not coincide
\item EU limited to coordinating public goods requiring collective action (defense, cross-border externalities)
\end{itemize}

This preserves member state sovereignty but requires accepting the costs of occasional sovereign defaults and fiscal crises at the member state level.

\subsubsection{Option C: Muddling Through (The Unstable Intermediate)}

Current arrangements approximate neither Option A nor Option B. The eurozone has:
\begin{itemize}
\item A common monetary policy but no fiscal union
\item De facto bailout mechanisms (ESM, OMT) but no explicit mutualization
\item Fiscal rules (Stability and Growth Pact, Fiscal Compact) that are not fully credible
\item No EU-level taxation authority
\item Banking union attempts without complete fiscal backing
\end{itemize}

The American historical experience suggests such intermediate arrangements are difficult to sustain. The United States spent less than a decade in a comparable intermediate state (Articles of Confederation, 1781--1788) before political crisis forced a choice. Europe has maintained its intermediate state much longer, perhaps because the costs have been distributed gradually across time (slow growth, high unemployment, periodic crises) rather than concentrated in a single dramatic crisis (like Shays' Rebellion).

\subsection{The Differences That Matter}

Several important differences between the 1780s United States and contemporary Europe affect the applicability of these lessons:

\begin{enumerate}
\item \textbf{Initial debt levels}: US debt in 1790 was about 40\% of GDP, incurred for a widely endorsed national purpose. European member state debts today are larger (averaging higher percentages of GDP) and were incurred for diverse national purposes.

\item \textbf{Political integration}: The American states in 1790 shared language, legal traditions, political culture, and recent common experience in the Revolution. European member states are more diverse culturally, linguistically, and institutionally.

\item \textbf{Labor and factor mobility}: Labor mobility among American states has historically been much higher than among European member states, providing an adjustment mechanism to asymmetric shocks.

\item \textbf{Size of government}: In the 1790s, federal government expenditures were about 1--2\% of GDP. European governments today are much larger (typically 40--50\% of GDP at all levels of government combined), making fiscal integration more consequential.

\item \textbf{Social insurance systems}: Modern European states have extensive pension, healthcare, and social insurance systems that differ substantially across countries. Fiscal integration would require either harmonizing these systems or allowing substantial transfers across countries with different social preferences.

\item \textbf{Democratic legitimacy}: The US Constitution was ratified through a deliberate constitutional process. European fiscal integration would require a similar democratic constitutional moment, which has not yet occurred and may not be politically feasible.
\end{enumerate}

These differences suggest that European fiscal integration would be more difficult and more consequential than American fiscal integration in 1790. The diversity of European member states and the size of their governments mean that fiscal union would require more substantial transfers and more difficult political compromises.

\subsection{A Conditional Recommendation}

Given this analysis, my conditional recommendation for Europe is:

\textbf{If European leaders want to preserve the eurozone}, they should move deliberately toward Option A (genuine fiscal union) following the 1790 precedent:
\begin{itemize}
\item Mutualize existing debts as part of a constitutional bargain
\item Simultaneously transfer taxation authority to EU level
\item Create strong EU-level fiscal capacity
\item Impose binding constraints on member state fiscal policies
\item Accept the political integration this requires
\end{itemize}

This would require a constitutional convention or comparable political process to establish democratic legitimacy. It would effectively create a federal European state, at least in fiscal matters.

\textbf{If European leaders are unwilling to accept such deep political integration}, they should instead move toward Option B (decentralized system):
\begin{itemize}
\item Strengthen no-bailout commitments credibly, accepting that this requires allowing defaults
\item Require member states to adopt constitutional fiscal constraints
\item Break bank-sovereign doom loops through regulatory means
\item Consider allowing exit from the eurozone for members whose fiscal and economic structures are incompatible with monetary union
\item Accept reduced fiscal ambitions at both EU and member state levels
\end{itemize}

This would preserve member state sovereignty but would require accepting greater economic volatility and occasional sovereign debt crises.

\textbf{What Europe should not do} is continue indefinitely in the current intermediate state. The American experience from the 1780s through the 1840s suggests that attempting to maintain common monetary institutions without clear resolution of fiscal authority creates persistent problems:
\begin{itemize}
\item Uncertainty about ultimate fiscal responsibilities impairs market functioning
\item Moral hazard problems accumulate as actors try to game ambiguous rules
\item Economic efficiency suffers from lack of clear institutional boundaries
\item Political legitimacy erodes as citizens cannot identify who is responsible for fiscal outcomes
\end{itemize}

\section{Conclusion}

Incorporating the symmetric competition between state and federal assumption plans in the 1780s fundamentally changes how we understand American fiscal union. Hamilton's 1790 assumption was not a simple bailout or an act of fiscal generosity. It was the winning move in a strategic competition to establish the federal-dominant equilibrium over the state-dominant equilibrium that many Americans preferred.

This revised understanding yields more nuanced lessons for contemporary debates about fiscal unions. Most importantly, it clarifies that federal assumption of lower-level government debts makes economic and political sense when it is part of a constitutional bargain that also transfers fiscal authority to the federal level. It does not make sense as an ongoing insurance arrangement where lower-level governments retain fiscal sovereignty.

For Europe, this implies a genuine choice: move toward deep fiscal integration (the 1790 path) or move toward a more decentralized confederal system with credible no-bailout commitments (the reformed Articles path). The current intermediate arrangement is likely not sustainable indefinitely.

The United States took about 60 years (from 1790 to 1850) to fully clarify the boundaries of federal fiscal authority, passing through one more major crisis (the 1840s state defaults) before the system stabilized. Europe has been in its intermediate state for about 25 years (since Maastricht in 1992). The American experience suggests that further clarification will come, either through deliberate constitutional reform or through crisis-driven adaptation. The question is whether European leaders will choose their path deliberately or allow events to choose for them.

\section*{References}

\begin{small}
\begin{itemize}
\item Adams, Henry C. 1887. \textit{Public Debts: An Essay in the Science of Finance}. New York: Appleton.

\item Ferguson, E. James. 1961. \textit{The Power of the Purse: A History of American Public Finance, 1776--1790}. Chapel Hill: University of North Carolina Press.

\item Kareken, John H., and Neil Wallace. 1978. ``Deposit Insurance and Bank Regulation: A Partial-Equilibrium Exposition.'' \textit{Journal of Business} 51 (3): 413--38.

\item Rothbard, Murray N. 1979. \textit{Conceived in Liberty, Volume 4: The Revolutionary War, 1775--1784}. New Rochelle, NY: Arlington House.

\item Additional references as in the original Nobel lecture.
\end{itemize}
\end{small}

\end{document}
